% ===========================================================================
% WAVVY REPORT TEMPLATE - v.1.0.0
% ===========================================================================

\documentclass[a4paper,11pt,table,dvipsnames]{article}

% ===========================================================================
% PACKAGES
% ===========================================================================

% Fonts and symbols
\usepackage{lmodern}
\usepackage{amsmath,amssymb}
\usepackage[utf8]{inputenc}
\usepackage[T1]{fontenc}
\usepackage{inconsolata}
\usepackage{wasysym}
\usepackage{microtype}

% Layout, page setup
\usepackage[margin=2.5cm,a4paper]{geometry}
\usepackage{titlesec}
\usepackage{appendix}
\usepackage{multicol}
\usepackage{multirow}
\usepackage{parskip}
\usepackage{tabto}
\usepackage{pdflscape}
\usepackage{enumitem}
\usepackage{adjustbox}
\usepackage[super]{nth}

% Science and language typesetting
\usepackage{bm}
\usepackage[
  arc-separator = \,,
  retain-explicit-plus,
  inter-unit-product = \cdot,
  detect-weight=true,
  detect-family=true,
  range-phrase=--,
  range-units=single
]{siunitx}
\usepackage[version=4]{mhchem}
\usepackage{nicematrix}
\usepackage[makeroom]{cancel}
\usepackage{circledsteps}
\usepackage[italian]{babel}
\usepackage{csquotes}

% Graphics, captions, tables
\usepackage{graphicx}
\usepackage{float}
\usepackage{tikz,tikz-3dplot}
\usepackage{pgfplots}
\usepackage{pdfpages}
\usepackage[justification=centering, labelfont={small,bf}, font={small}]{caption}
\usepackage{subcaption}
\usepackage{array}
\usepackage{tabularx}
\usepackage{booktabs}
\usepackage[outline]{contour}

% TikZ options setup
\usetikzlibrary{
  calc,
  arrows,
  arrows.meta,
  positioning,
  decorations.pathreplacing,
  decorations.markings,
  decorations.text,
  calligraphy,
  pgfplots.dateplot,
  shapes.geometric
}

% References and links
\usepackage{hyperref}
\usepackage[noabbrev]{cleveref}
\usepackage[nottoc,numbib]{tocbibind}

% Miscellaneous
\usepackage{datetime2}
\usepackage{iftex,etoolbox,expl3,xparse}

\pgfplotsset{compat=1.18}
\AtBeginDocument{\RenewCommandCopy\qty\SI}

% ===========================================================================
% METADATA
% ===========================================================================

\title{\textbf{[Nome Feature]}}
\author{\textbf{[Nome Autore]} -- Wavvy}
\date{Versione [X.Y.Z] \\ \today}

\begin{document}

\maketitle

\tableofcontents

\newpage

% ===========================================================================
% TEMPLATE CONTENT
% ===========================================================================

\section{Templates (da rimuovere)}

% Table template
\subsection{Tabelle e diagrammi}

\subsubsection{Tabella}
% Table === BEGIN
\begin{table}[H]
    \centering
    \renewcommand{\arraystretch}{1.2}
    \begin{tabular}{@{} l l c @{}}
        \toprule
        \textbf{Titolo 1} & \textbf{Titolo 2} & \textbf{Titolo 3} \\
        \midrule
        Riga 1            & Riga 1            & Riga 1            \\
        Riga 2            & Riga 2            & Riga 2            \\
        Riga 3            & Riga 3            & Riga 3            \\
        \midrule
        \textbf{Titolo 1} &                   & Riga 4            \\
        Riga 5            &                   & Riga 5            \\
        \bottomrule
    \end{tabular}
    \caption{Descrizione della tabella}
    \label{tab:nome_tabella}
\end{table}
% Table === END

\subsubsection{Diagramma}
\begin{figure}[H]
    \centering
    \begin{tikzpicture}[
            node distance=2cm,
            auto,
            font=\small\sffamily,
            % Stili dei nodi
            startstop/.style={rectangle, rounded corners, minimum width=3cm, minimum height=1cm, text centered, draw=black, fill=gray!10},
            process/.style={rectangle, minimum width=3cm, minimum height=1cm, text centered, draw=black, fill=white},
            decision/.style={diamond, aspect=2, minimum width=3cm, minimum height=1cm, text centered, draw=black, fill=gray!5},
            arrow/.style={thick,->,>=stealth}
        ]

        % Nodi
        \node (start) [startstop] {Evento Iniziale};
        \node (check) [decision, below of=start, yshift=-0.5cm] {Decisione};

        % Diramazioni
        \node (fail) [process, below of=check, xshift=-3.5cm, yshift=-0.5cm] {Evento};
        \node (success) [process, below of=check, xshift=3.5cm, yshift=-0.5cm] {Evento};

        % End
        \node (end) [startstop, below of=check, yshift=-3cm] {Evento Finale};

        % Frecce
        \draw [arrow] (start) -- (check);
        \draw [arrow] (check) -| node[anchor=south east] {Condizione} (fail);
        \draw [arrow] (check) -| node[anchor=south west] {Condizione} (success);
        \draw [arrow] (fail) |- (end);
        \draw [arrow] (success) |- (end);

    \end{tikzpicture}
    \caption{Nome del diagramma di flusso}
    \label{fig:nome_diagramma}
\end{figure}


\subsection{Formule matematiche}

\subsubsection{Definizione}
% Definition === BEGIN
\textbf{Definizione (\(NomeParametro\)):}
Sia
\[
    Parametro = \dots
\]
\dots .
% Definition === END

\subsubsection{Algoritmo}
% Algorithm === BEGIN
\textbf{Algoritmo (\(NomeAlgoritmo\)):}
Sia
\[
    Algoritmo = \dots
\]
la algoritmo che \dots . L'algoritmo prevede:
\begin{enumerate}
    \item Passo 1
    \item Passo 2
    \item Passo 3
\end{enumerate}
% Algorithm === END

\subsubsection{Regola}
% Rule === BEGIN
\textbf{Regola:}
\dots
% Rule === END

% ===========================================================================
% BODY
% ===========================================================================

\section{Scopo e perimetro}

Questo documento descrive l'implementazione e le scelte architetturali relative a \textbf{[Nome Feature]}.

\begin{itemize}
    \item Piattaforme coinvolte: \dots
    \item Tecnologie: \dots
\end{itemize}

\subsection{Perimetro}
Questo documento si concentra sui seguenti aspetti:
\begin{itemize}
    \item Punto 1
    \item Punto 2
\end{itemize}

Sono esplicitamente esclusi:
\begin{itemize}
    \item Punto 1
\end{itemize}

\section{Panoramica architetturale}
L'architettura del [Nome Feature] \dots

\section{Principi di progettazione}
Il [Nome feature] segue i seguenti principi di progettazione:

\subsection{Principio 1}
\dots

\section{(Spiegazione della feature)}
\dots

% Keep everything below this line at the end of the document
% ===========================================================================
\section{Flusso \dots}
\dots

\section{Threat model e mitigazioni}
\dots

\section{Alternative considerate e modelli mitigati}
\dots

\section{Sintesi}
\dots



\end{document}